\documentclass[a4paper, serif, 10pt]{moderncv}

%% moderncv
% style and color
\moderncvstyle{banking}
\moderncvcolor{black}

%% page layout/margins
\usepackage[top=2.5cm, bottom=2.5cm, left=1.5cm, right=1.5cm]{geometry}

\nopagenumbers{}

%% Babel config for German language
% ref:
% - https://latex3.github.io/babel/guides/locale-german.html
\usepackage[ngerman]{babel}

%% fontspec
\usepackage{fontspec}

\IfFontExistsTF{Linux Libertine}{
  \setmainfont[Ligatures={TeX, Common}, Numbers=OldStyle]{Linux Libertine}
}{}
\IfFontExistsTF{Linux Libertine O}{
  \setmainfont[Ligatures={TeX, Common}, Numbers=OldStyle]{Linux Libertine O}
}{}

%% CTeX
% CJK support, used to show CJK characters in the resume
%
% - fontset=none: disable builtin fontset but instead set the CJK font manually
% - heading=false: disable ctex heading
% - punct=kaiming: use kaiming punctuations styles for CJK
% - scheme=plain: use plain scheme, do not override \normalsize font size
% - space=auto: space settings for CJK characters
%
% ref:
% - http://ctan.mirrorcatalogs.com/language/chinese/ctex/ctex.pdf
\usepackage[UTF8, heading=false, punct=kaiming, scheme=plain, space=auto]{ctex}

\IfFontExistsTF{Noto Serif CJK SC}{
  \setCJKmainfont{Noto Serif CJK SC}
}{}
\IfFontExistsTF{Noto Sans CJK SC}{
  \setCJKsansfont{Noto Sans CJK SC}
}{}

%% line spacing
\usepackage{setspace}
\setstretch{1.125}

%% URL styling - use normal text instead of monospace
\urlstyle{same}

% Auto-underline all links
% Defer until \begin{document} so hyperref is already loaded
\AtBeginDocument{%
  \let\oldhref\href
  \renewcommand{\href}[2]{\oldhref{#1}{\underline{#2}}}%
}

%% Patch \httplink and \httpslink to handle full URLs
% The original moderncv macros always prepend http:// or https:// to URLs.
% This causes issues when the URL already has a protocol (e.g., https://example.com)
% resulting in malformed URLs like https://https://example.com.
% This patch checks if the URL already contains "://" and uses it directly if so.
% Note: We use \str_set:Nx (x-expansion) to expand macros like \@homepage first.
\ExplSyntaxOn
\str_new:N \g_yamlresume_url_str

\RenewDocumentCommand{\httpslink}{O{}m}{
  \str_set:Nx \g_yamlresume_url_str {#2}
  \str_if_in:NnTF \g_yamlresume_url_str {://}
    { \url{#2} }
    { \url{https://#2} }
}

\RenewDocumentCommand{\httplink}{O{}m}{
  \str_set:Nx \g_yamlresume_url_str {#2}
  \str_if_in:NnTF \g_yamlresume_url_str {://}
    { \url{#2} }
    { \url{http://#2} }
}
\ExplSyntaxOff

\name{Katrin Vogel}{}
\title{UX Designerin mit Fokus auf nutzerzentrierte Gestaltung}
\phone[mobile]{+49 151 23456789}
\email{hallo@katrinvogel.de}
\homepage{https://katrinvogel.de}

\address{München, Bayern, Deutschland, 80331}{}{}

\extrainfo{{\small \faLinkedin}\ \href{https://www.linkedin.com/in/katrinvogel}{@katrinvogel} {} {} {} • {} {} {}
{\small \faBehance}\ \href{https://www.behance.net/katrinvogel}{@katrinvogel} {} {} {} • {} {} {}
{\small \faDribbble}\ \href{https://dribbble.com/katrinvogel}{@katrinvogel}}

\begin{document}

\maketitle

\section{Basisinformationen}

\cvline{}{Leidenschaftliche UX Designerin mit über 6 Jahren Erfahrung in der Gestaltung digitaler Produkte, die Menschen begeistern und im Alltag echten Mehrwert schaffen. Ich verbinde tiefes User Research mit klarem visuellem Design, um barrierefreie, intuitive und ästhetisch überzeugende Erlebnisse zu entwickeln. Meine Expertise umfasst Design Thinking, Prototyping, Design Systems sowie die enge Zusammenarbeit mit Product Management und Entwicklung.
%
\begin{itemize}
\item Schwerpunkt auf nutzerzentrierten Prozessen und datenbasierten Designentscheidungen
\item Erfahrung in der Leitung von Design Sprints und Workshops
\item Starkes Gespür für Typografie, Farbe und visuelle Hierarchie
\item Überzeugte Verfechterin von Barrierefreiheit und Inklusion im digitalen Raum
\end{itemize}}
\section{Ausbildung}

\cventry{Okt. 2015–Sept. 2019}
        {Bachelor, Kommunikationsdesign, Punkte: 1,6}
        {Hochschule für Gestaltung München}
        {\url{https://www.hfg-muenchen.de}}
        {}
        {\begin{itemize}
\item Vertiefung in Interaktions- und Interface-Design mit Fokus auf digitale Produkte
\item Praxisprojekte in Kooperation mit Unternehmen aus den Bereichen FinTech und Mobilität
\item Abschlussarbeit zum Thema „Inclusive Design – Barrierefreiheit in digitalen Produkten“ mit der Note 1,3
\item Studentische Hilfskraft im UX-Labor, Betreuung von Usability-Studien
\end{itemize}
\textbf{Kurse}: Grundlagen der Gestaltung, Typografie, Interaktionsdesign, Usability Engineering, Design Research, Bewegtbild und Animation}
\section{Berufserfahrung}

\cventry{März 2022–Heute}
        {Senior UX Designerin}
        {Pixelwerk GmbH}
        {\url{https://pixelwerk.de}}
        {}
        {\begin{itemize}
\item Verantwortlich für die UX-Strategie und die Weiterentwicklung des Design Systems der SaaS-Plattform „DataFlow“
\item Leitung von User-Research-Studien mit über 40 qualitativen Interviews und iterativen Usability-Tests
\item Etablierung eines Design-Sprint-Formats, das die Time-to-Market neuer Features um 30 \% reduzierte
\item Enge Zusammenarbeit mit Product Management, Engineering und Data Science zur Integration nutzerzentrierter Prozesse
\item Mentoring von Junior UX Designerinnen und Designern sowie Aufbau einer internen Wissensdatenbank
\end{itemize}
\textbf{Schlüsselwörter}: Design Thinking, Design System, User Research, Design Sprints, Mentoring}

\cventry{Sept. 2019–Feb. 2022}
        {UX Designerin}
        {Appsphere Solutions}
        {\url{https://appsphere-solutions.de}}
        {}
        {\begin{itemize}
\item Konzeption und Gestaltung von User Interfaces für iOS- und Android-Apps im Bereich Health und E-Commerce
\item Durchführung von Stakeholder-Workshops und Ableitung von UX-Strategien für neue Produktfunktionen
\item Erstellung von interaktiven Prototypen mit Figma und InVision zur frühzeitigen Validierung von Ideen
\item Durchführung von A/B-Tests und Ableitung datenbasierter Designentscheidungen
\item Aufbau eines nutzerzentrierten Designprozesses in einem agilen Entwicklungsteam
\end{itemize}
\textbf{Schlüsselwörter}: Mobile Apps, Prototyping, Usability Testing, A/B-Tests, Agile}
\section{Sprachen}

\cvline{Deutsch}{Muttersprache oder zweisprachig \hfill \textbf{Schlüsselwörter}: Muttersprache}
\cvline{Englisch}{Verhandlungssichere Kenntnisse \hfill \textbf{Schlüsselwörter}: Business English, UX-Workshops}
\section{Fähigkeiten}

\cvline{UX Research}{Erfahren \hfill \textbf{Schlüsselwörter}: Tiefeninterviews, Usability-Tests, Personas, Customer Journey, Umfragen}
\cvline{UI-Design}{Experte \hfill \textbf{Schlüsselwörter}: Figma, Sketch, Adobe XD, Design Systems, Barrierefreiheit}
\cvline{Prototyping}{Erfahren \hfill \textbf{Schlüsselwörter}: Figma Prototype, InVision, ProtoPie, Micro-Interactions}
\cvline{Methoden}{Erfahren \hfill \textbf{Schlüsselwörter}: Design Thinking, Design Sprints, Service Blueprint, Jobs-to-be-Done}
\section{Auszeichnungen}

\cventry{Feb. 2023}
        {Rat für Formgebung}
        {German Design Award 2023}
        {}
        {}
        {Auszeichnung für herausragendes Interface-Design der mobilen App „FairTrade Finder“ in der Kategorie „Excellent Communications Design“.}
\section{Zertifikate}

\cventry{Jan. 2021}
        {Google}
        {Google UX Design Professional Certificate}
        {\url{https://www.coursera.org/professional-certificates/google-ux-design}}
        {}
        {}

\cventry{Nov. 2020}
        {Fraunhofer-Institut für Arbeitswirtschaft und Organisation}
        {Certified Professional for Usability and User Experience}
        {\url{https://www.iao.fraunhofer.de}}
        {}
        {}
\section{Veröffentlichungen}

\cventry{März 2023}
        {Barrierefreiheit im digitalen Produktdesign – Ein Leitfaden für Teams}
        {UX-Love Magazin}
        {\url{https://ux-love-magazin.de/artikel/barrierefreiheit}}
        {}
        {\begin{itemize}
\item Praxisorientierter Leitfaden zur Integration von Barrierefreiheit in agile Designprozesse
\item Vorstellung von Methoden und Tools für inklusive Nutzertests
\item Plädoyer für Barrierefreiheit als Gestaltungsprinzip statt nachträglicher Anforderung
\end{itemize}}
\section{Projekte}

\cventry{Sept. 2021–Feb. 2022}
        {Mobiler Produktbegleiter für nachhaltigen Konsum}
        {FairTrade Finder}
        {\url{https://fairtrade-finder.de}}
        {}
        {\begin{itemize}
\item Konzeption einer App, die Verbraucher:innen transparente Informationen zu Herkunft und Nachhaltigkeit von Produkten bietet
\item Durchführung von User Research mit 25 Teilnehmenden zur Identifikation von Entscheidungsmustern im Kaufprozess
\item Erstellung von interaktiven Prototypen und iterative Usability-Tests über drei Zyklen
\item Finale Gestaltung des UI-Kits unter Berücksichtigung der WCAG-2.1-Richtlinien
\end{itemize}
\textbf{Schlüsselwörter}: Mobile App, Nachhaltigkeit, User Research, WCAG}

\cventry{Jan. 2022–Dez. 2023}
        {Open-Source Design System für barrierefreie Webanwendungen}
        {Design System „Kiosk“}
        {\url{https://github.com/katrinvogel/kiosk}}
        {}
        {\begin{itemize}
\item Aufbau eines modular aufgebauten Design Systems mit über 60 Komponenten in Figma
\item Entwicklung von Dokumentationsstandards für Design-Tokens und Nutzungskontexte
\item Aktive Pflege der Community und Organisation von monatlichen Contributing-Workshops
\end{itemize}
\textbf{Schlüsselwörter}: Design System, Figma, Open Source, Accessibility}
\section{Interessen}

\cvline{Design}{Design Thinking, Generative KI, Nachhaltiges Design}
\cvline{Musik}{Klavier, Konzerte}
\cvline{Reisen}{Skandinavien, Städte mit Designkultur}
\section{Ehrenamtliche Tätigkeiten}

\cventry{Sept. 2022–Heute}
        {Mentorin für UX Design}
        {ReDI School of Digital Integration}
        {\url{https://www.redi-school.org}}
        {}
        {\begin{itemize}
\item Unterstützung geflüchteter und zugewanderter Menschen beim Einstieg in den UX-Design-Bereich
\item Planung und Durchführung von monatlichen Feedback-Sessions zu Portfolio-Projekten
\item Aufbau von Lernmaterialien und Übungen zu den Grundlagen nutzerzentrierter Gestaltung
\end{itemize}}

\end{document}